\documentclass[9pt, a4paper]{article}
% ======================
% 基础文档设置与中文支持
% ======================
\usepackage{geometry}           % 页面布局控制（边距、纸张大小等）
\geometry{scale=0.75}
\usepackage{setspace}           % 行间距控制
\usepackage{fontspec}           % 高级字体控制
\usepackage{xeCJK}              % 中文字体支持
\usepackage{ctex}               % 完整中文支持
\usepackage{CJKutf8}            % UTF-8编码支持（备用）
\usepackage{zhnumber}           % 中文数字转换

% ======================
% 数学相关包
% ======================
\usepackage{amsmath}            % 基础数学公式（矩阵、多行公式等）
\usepackage{amssymb}            % 扩展数学符号
\usepackage{amsthm}             % 定理环境定制
\usepackage{mathtools}          % amsmath的增强版（矩阵、符号等）
\usepackage{mathrsfs}           % 提供花体字母
\usepackage{latexsym}           % 特殊数学符号
\usepackage{stmaryrd}            % 额外数学符号

% ======================
% 图形与绘图
% ======================
\usepackage{graphicx}           % 图片插入（\includegraphics）
\usepackage{tikz}               % 矢量绘图核心包
\usepackage{tikz-cd}            % 交换图绘制（范畴论）
%\usepackage{quiver}             

% ======================
% 参考文献与引用
% ======================
\usepackage[authoryear]{natbib} % 作者-年份引用格式
\usepackage{etoolbox}           % 参考文献工具（定制bib样式）

% ======================
% 超链接与交互
% ======================
\usepackage{xcolor}             % 颜色支持（链接着色）
\usepackage{hyperref}           % 超链接/书签（需最后加载）
\hypersetup{                    % hyperref配置
    colorlinks=true,            % 彩色链接而非方框
    citecolor=green,            % 引用颜色
    urlcolor=red,               % 外部URL颜色
    bookmarksnumbered=true      % 书签显示章节编号
}
\usepackage{cleveref}           % 智能引用（\cref自动识别类型）

% ======================
% 列表与布局增强
% ======================
\usepackage{enumitem}           % 列表定制（间距、标签等）
\usepackage{fancyhdr}           % 页眉页脚定制
\usepackage{lmodern}            % 拉丁现代字体（矢量字体）

% ======================
% 其他工具
% ======================
\usepackage{pdfpages} 

% 定理环境样式
\theoremstyle{plain}
\newtheorem{theorem}{定理}
\newtheorem{lemma}[theorem]{引理}
\newtheorem{corollary}[theorem]{推论}
\newtheorem{proposition}[theorem]{命题}

\theoremstyle{definition}
\newtheorem{definition}[theorem]{定义}
\newtheorem{example}[theorem]{示例}
\newtheorem{remark}[theorem]{注记}

\theoremstyle{remark}

% 自定义命令
\definecolor{arxivblue}{RGB}{0, 102, 204}
\newcommand{\arxiv}[3]{\href{https://arxiv.org/abs/#1}{#1}, \textbf{\textcolor{arxivblue}{#2}}, #3}
\newcommand{\AGnumber}{0}
\newcommand{\RTQAnumber}{22}
\newcommand{\NewestDate}{2026年11月11日}

\newcommand{\todayarXiver}{} %今天的整理者！



\begin{document}
\begin{center}
    \Huge\textbf{不漏arXiv：\NewestDate\footnote{感谢arXiv提供的服务.}\footnote{今天的整理者是\todayarXiver}\footnote{由于arXiv api的局限性, 实际上获取的是arXiv api给出的最近三天submit的所有论文, 和recent页面可能有一定的差异. 三个category中没有进行去重, 顺序也可能存在一定问题, 麻烦整理者自行增删.}} \\
\end{center}

\setcounter{section}{-1}

\begin{center}
    今日AG数量：\AGnumber
    
    今日RT \& QA数量：\RTQAnumber
    
\end{center}


\section{精选}


\section{AG}
%AG begin
%AG end

\section{RT \& QA}
%RT&QA begin
\arxiv{2611.14304v1}{Noetherian connected braided Hopf algebras are affine}{Huan Jia, Yinhuo Zhang}

\arxiv{2611.14084v1}{Representations of the quantum double of the restricted quantum group $\mathbf{\overline{u}}_q(\mathfrak{sl_2})$}{Hua Sun, Huixiang Chen, Xiaoming Li 等}

\arxiv{2611.13974v1}{Two-parameter versus one-parameter quantum groups: structural differences arising from the type~$D$ phenomenon}{Naihong Hu, Marc Rosso, Xiao Xu}

\arxiv{2611.13897v1}{On the strong Dixmier-Moeglin equivalence for the quantum semisimple groups}{Nirina Albert Razafimandimby}

\arxiv{2611.13803v1}{Elliptic hypergeometric sum/integrals for orthogonal groups and Seiberg duality on the lens space}{Andrew P. Kels, Masahito Yamazaki}

\arxiv{2611.13639v1}{Quantum Langlands correspondence for Painlevé II}{Sibasish Banerjee, Babak Haghighat, Jie Sun}

\arxiv{2611.13446v1}{Differential Monodromy, $q$-Opers, and Universal $(G,q)$-Wronskians}{Anton M. Zeitlin}

\arxiv{2611.13424v1}{Ostrik's Theorem for Operator Algebras in a Unitary Tensor Category}{Emeline Sherman}

\arxiv{2611.13347v1}{Computational Power of Clifford and Ising Quantum Computers}{Sawyer Dobson, Zhenghan Wang}

\arxiv{2611.13239v1}{A decomposition theorem for minimal W-algebras of type A}{Thibault Juillard}

\arxiv{2611.13129v1}{The uncertainty operator, the Weil ground state and the zeros of zeta: a numerical study and conditional positivity of the Weil quadratic form}{Ronnie Andrews, Alain Connes, Antonio La Gatta}

\arxiv{2611.12874v1}{FRT presentation of the $q$-Onsager comodule algebra and refined $q$-characters}{Pascal Baseilhac, Azat M. Gainutdinov}

\arxiv{2611.12836v1}{Property F implies weak integrality for non-degenerate pivotal braided fusion categories}{Justin Asher, Jonah Berean-Dutcher, Clay Cordova 等}

\arxiv{2611.12723v1}{Classification of quasi-lisse simple affine vertex algebras}{Drazen Adamovic, Xuanzhong Dai, Kaiwen Sun}

\arxiv{2611.12417v1}{Holomorphic integer-graded vertex operator superalgebras are purely even}{Chongying Dong, Kaiwen Sun, Haowu Wang}

\arxiv{2611.12191v1}{The multiplicity bound for near-group fusion categories of type $\mathbb Z/4\mathbb Z+k$}{Yue Meng, Zhiqiang Yu}

\arxiv{2611.11699v1}{Frobenius and ozone automorphism groups of root of unity quantum cluster algebras}{Yunmeng Wu, Milen Yakimov}

\arxiv{2611.11685v1}{Rational knots, quivers, and tails of colored HOMFLY-PT polynomials}{Elijah Berry, Rongjiang Guo, Jonathan A. Higgins 等}

\arxiv{2611.11672v1}{Pointed Hopf Algebras from Pale Braidings in Characteristic 2}{Dirceu Bagio, Laura Natalia Orozco García}

\arxiv{2611.11575v1}{Representations of the Weyl algebra modulo $p^2$}{Ruben Mamani-Velasco, Akaki Tikaradze}

\arxiv{2611.11502v1}{Does Merge Form a Hopf Algebra? Coproducts, Traces, and the Prospects of Algebraic Minimalism}{Daisuke Bekki}

\arxiv{2611.11240v1}{Stretched near-groups}{Andrew Schopieray, Yilong Wang, Zhiqiang Yu}

%RT&QA end


\end{document}